%% Relative references under beamer: the class that anchors nothing for us.
%%
%% beamer sets hyperref's implicit=false -- it anchors its own \labels and
%% runs its own navigation -- so there is no destination at
%% \refstepcounter to link to.  linguexx used to read that as "links are
%% impossible here" and switch itself off, which left \ref moving and
%% \Last not: in the class most linguistics slides are written in, exactly
%% the inconsistency the links exist to remove.  Nothing about
%% implicit=false makes a destination impossible, so linguexx now makes
%% one itself.
%%
%% The second half of the case is the overlays, and it is the half that
%% cannot be seen on paper.  A frame is typeset once per slide with the
%% example counters restored each time, so every example on the two-slide
%% frame below comes past TWICE with the same number.  Unguarded that is a
%% duplicate destination (which hyperref drops, warning once per repeat)
%% and a duplicate .aux record -- and the second record would be read as
%% two examples claiming one anchor, so the guard against SHARED anchors
%% would then refuse to link the very examples this case is about.  One
%% example seen twice has to be told from two examples numbered alike,
%% and the .aux is where that shows.
%% beamer, the class this package's users give talks in -- and the reason
%% the relative-reference links needed a second way of anchoring an
%% example.  beamer loads hyperref with implicit=false, so hyperref places
%% no destination at \refstepcounter and linguexx places its own.
%%
%% No theme and no navigation symbols: the case is about destinations, and
%% a theme would put dozens of its own links on every slide.
\documentclass{beamer}
\usetheme{default}
\setbeamertemplate{navigation symbols}{}
\usepackage{iftex}
\ifPDFTeX
  \usepackage[T1]{fontenc}
  \usepackage[utf8]{inputenc}
\fi
\usepackage{linguexx}

\begin{document}
\begin{frame}{BMFRAMEONE}
\ex.\label{bm:one} BMONE first example.

BMREF \ref{bm:one} BMLAST \Last\ BMNEXT \Next
\end{frame}

%% Two slides, two passes over everything in here.
\begin{frame}{BMFRAMETWO}
\ex. BMTWO second example.
\only<2->{BMOVERLAY only on the second slide.}

BMLASTTWO \Last
\end{frame}

%% \pause: an example EXECUTED on the frame's first slide and visible only
%% on the next.  This is the commoner half of the problem and the one that
%% reached a reader.  \pause, \uncover and \onslide run their material on
%% every slide of the frame -- the counter steps, and an anchor placed
%% there is placed on that slide -- and drop its ink on the slides where it
%% is covered.  Anchoring the example where it is first run therefore sends
%% the click to a slide where the example cannot be seen, and nothing on
%% the page shows it: there is nothing on that page to show.
%%
%% The two examples below are on one frame, one slide apart, and that is
%% the whole assertion: their destinations must be one page apart too.
%% Anchored where they are first run, both would land on the same page.
\begin{frame}{BMFRAMETHREE}
\ex. BMSHOWN visible at once.
\pause
\ex. BMPAUSED typeset here, visible only on the next slide.
\end{frame}

%% An example whose FIRST appearance is on a later slide.  \only does not
%% typeset its argument at all on the slides it excludes, so this example
%% does not exist on slide 1 of this frame and is numbered for the first
%% time on slide 2.  A rule that asked "is this the frame's first pass"
%% skipped it outright and gave it no anchor, so BMBEFORE below -- which
%% names it -- was reported as pointing at an example that does not exist,
%% by the mechanism that had just discarded its target.  Asking instead
%% whether the name is new TO THIS FRAME gets it right.
%%
%% Last in the document on purpose: with \only around a numbered example
%% the count itself differs between the slides of this frame, so anything
%% numbered after it would differ too.
\begin{frame}{BMFRAMEFOUR}
BMBEFORE \Next

\only<2->{\ex. BMLATE only on the second slide.}
\end{frame}
\end{document}
